\documentclass{article}
\usepackage{/globals/de}
\usepackage{/globals/anki}
\ankiprefix{systemintegration}
\ankitopic{Systemintegration}
\ankishow
\begin{document}
\section{Integration}
Ein System kann auf mehrere Art und Weisen integriert sein.
\begin{description}
 \item[Datenintegration] Auch wenn ein Datensatz mehrfach und andersweitig genutzt wird, wird er nur ein mal gespeichert. So kann es nicht zu inkonsistenz kommen.
 \item[Funktionsintegration] Ein Anwendungssystem hat mehrere Funktionen, auch Aktivitäten genannt. Dessen Modelle sind \idr nach den Bereichen der Organisation des Unternehmens sortiert.
 \item[Prozessintegration] Alle Funktionen eines Prozesses sind verbunden.
\end{description}
\ankinote{arten}{Arten}{Daten, Funktions, Prozess}
Informationen können sowohl horizontal als auch vertikal integriert sein.
\begin{description}
 \item[Horizontal] Die Administrations- und Dispositionssysteme unterschiedliche Unternehmensbereiche sind verbunden.
 \item[Vertikal] Die Planungs- und Kontrollsysteme sind mit den Daten der Administrations und Dispositionssysteme verbunden.
\ankinote{richtungen}{Richtungen}{Horizontal (AD von Bereichen verbunden), Vertikal (PK mit AD verbunden)}
\end{description}
Integrationen können desweiteren in ihrer Reichweite (Innerbetrieblich oder Zwischenbetrieblich) und ihrem Automatisierungsgrad (Vollautomatisiert oder Teilautomatisiert) unterschieden werden.
\ankinote{unt-arten}{Unterschied Arten}{Reichweite (Innerbetrieblich vs Zwischenbetrieblich), Automatisierungsgrad (Vollautomatisiert vs Teilautomatisiert)}

\end{document}